\begin{figure*}[!t]
\centering\color{revisedred}
\begin{tikzpicture}[
 constructionbox/.style={draw=teal!70,fill=wash,rounded corners=2pt,align=center,text width=4.7cm,minimum height=1.15cm,inner sep=5pt,font=\small},
 wide/.style={draw=teal!70,fill=white,rounded corners=2pt,align=center,text width=16.1cm,minimum height=0.8cm,inner sep=5pt,font=\small},
 connection/.style={-{Latex[length=2mm]},thick,draw=teal}]
\node[wide] (question) at (5.7,0) {Decision question, reference world, matched branches and observation horizon $T$};
\node[constructionbox] (actors) at (0,-1.7) {\textit{Agent swarm} and incentives\\Agent roles, resources, preference\\weights and feasible actions};
\node[constructionbox] (information) at (5.7,-1.7) {\textit{Social graph} environment\\Observation boundaries, directed links\\and swarm message delivery};
\node[constructionbox] (rules) at (11.4,-1.7) {Rules and uncertainty\\Settlement ledgers, event timing\\and declared assumptions};
\draw[connection] (question.south -| actors.north)--(actors.north);
\draw[connection] (question.south)--(information.north);
\draw[connection] (question.south -| rules.north)--(rules.north);
\node[wide] (review) at (5.7,-3.5) {Construction review\\Evidence and rationale; resource consistency; social-graph and interaction checks};
\draw[connection] (actors.south)--(actors.south |- review.north);
\draw[connection] (information.south)--(review.north);
\draw[connection] (rules.south)--(rules.south |- review.north);
\node[wide] (output) at (5.7,-5) {Reviewed world and \textit{agent swarm} $\rightarrow$ local agent decision model (\textit{Laya}) and execution ledgers};
\draw[connection] (review.south)--node[right,font=\footnotesize]{accepted after review}(output.north);
\draw[connection,dashed] (review.east)--++(0.55,0)|-node[pos=0.25,right,font=\footnotesize]{revise}(question.east);
\end{tikzpicture}
\caption{\rev{World and \textit{agent swarm} construction in \textit{Boyun},
inspired by \textit{MiroFish}'s world-building workflow.} A decision question
determines three complementary specification tasks \rev{defining the swarm of
agents, their social graph, and the settlement rules}. Their outputs undergo
structural checks and human review, with author-directed revision where needed.
The process specifies \rev{agent preference weights, social-graph links, and
settlement rules before \textit{Laya} selects live agent actions locally.}}
\label{fig:construction}
\end{figure*}

\section{Problem Formulation and World Creation}\label{sec:world-creation}
The preceding section identified the conceptual foundations needed to study
strategic interactions. To turn those foundations into a reproducible
experiment, \rev{\textit{Boyun}---following \textit{MiroFish}'s world-first
workflow---}constructs a \textit{world}: an explicit specification of \rev{an
\textit{agent swarm} (modeled actors operating as autonomous agents)},
resources, \rev{a \textit{social graph} governing observations and communication},
actions, events, and settlement rules. Construction is performed by a user or a
separate authoring model, followed by structural checks and human review, and
finishes before the \rev{agent decision model} is ever called.
Figure~\ref{fig:construction} summarizes this preparation; the execution
procedure follows in Section~\ref{sec:method-system}.

\subsection{Core Terminology}
Before specifying a world, a reader must distinguish the static specification
from the record produced by executing it. Table~\ref{tab:definitions} defines
worlds and their execution, while actor and decision terms appear in the
subsections below. In particular, a world is a model of a decision setting,
whereas a run is one execution of that model; two runs of the same world can
differ whenever the selection rule uses sampling. For an everyday
analogy, a board game's rules and starting setup define the world, whereas one
played game is a run and its recorded move sequence is a trajectory.

\begin{table}[H]
\caption{Worlds and execution terminology.}
\label{tab:definitions}
\centering\small
\renewcommand{\arraystretch}{1.12}
\begin{tabularx}{\columnwidth}{@{}>{\raggedright\arraybackslash}p{0.25\columnwidth}Y@{}}
\toprule
Term & Definition \\
\midrule
\textit{World} & A specified decision setting: actors, resources, observations, actions, events, and settlement rules. \\
\textit{Scenario} & An alternative setting examined through a world or matched branch; its label does not establish a probability. \\
\textit{State} & Ledger-specific data at step $t$: balances and metrics, plus reserves, collateral, debt, or claims where modeled. \\
\textit{Run / trajectory} & A run executes one world. Its trajectory records observations, actions, and settled states. \\
\textit{Round / horizon} & Index $t$ orders rounds or decision points up to $T$; elapsed time is recorded separately. A cycle is not inherently an hour or day. \\
\textit{Outcome / metric} & An outcome is a simulated result. A metric $x_{k,t}$ describes shared state within bounds $[L_k,H_k]$; metric points are not money. \\
\bottomrule
\end{tabularx}
\end{table}

Once those terms are separated, the author can state what the simulation
is designed to compare. Consequently, world creation begins with the decision
question rather than with narrative persona writing.

\subsection{Decision Question, Alternatives, and Horizon}
Specifically, the author first defines the focal decision, the alternatives
under comparison, the observation horizon $T$, and the outcomes of interest.
For example, a protocol designer might compare early reserve disclosure against
delayed disclosure and track partner participation alongside retained cash.
That question dictates which actors and ledgers belong in the world; adding
unrelated detail does not strengthen the comparison.

From that focal question, the author defines a reference world and
isolates which inputs change in each matched branch. Comparing disclosure
policies, for instance, keeps initial balances and unrelated rules fixed
wherever possible; if several inputs change together, the resulting difference
reflects their joint intervention rather than a single cause. Likewise,
the horizon $T$ bounds only the observation window: reaching $T$ does not imply
that the modeled protocol or market has ended.

\subsection{\rev{Agent Swarm}, Preference Weights, and Strategic Dependence}
With the comparison fixed, \rev{agents enter the \textit{swarm} (the world's
cast)} only if they control relevant resources, hold distinct information, or
can alter an outcome that another \rev{agent} values. An \rev{\textit{actor} in
\textit{Boyun} is an autonomous modeled \textit{agent}}, while a
\textit{persona} describes its role and motivation. To make that persona
executable \rev{within the \textit{swarm of agents}}, the world assigns explicit
goals, initial balances, feasible actions, \rev{social-graph} observation
access, and numeric preference weights, giving the \rev{agent swarm} a
game-theoretic structure rather than a list of biographies
\cite{polak2,nash1950}. Table~\ref{tab:actor-definitions} separates these
elements.

\begin{table}[H]
\caption{\rev{Agent swarm, goals, and incentives.}}
\label{tab:actor-definitions}
\centering\small
\renewcommand{\arraystretch}{1.12}
\begin{tabularx}{\columnwidth}{@{}>{\raggedright\arraybackslash}p{0.25\columnwidth}Y@{}}
\toprule
Term & Definition \\
\midrule
\rev{\textit{Actor / agent}} & \rev{An autonomous modeled decision-maker indexed by $i$ (a player in game theory); the interacting collection of agents forms the \textit{agent swarm} (or \textit{cast}).} \\
\textit{Persona} & The \rev{agent's} descriptive role and motivation, paired with explicit preference weights and actions. \\
\textit{Goal} & The outcome the \rev{agent} seeks, distinct from the analyst's research question. \\
\textit{Incentive} & A cost, benefit, or conditional effect that makes an action more or less attractive under the \rev{agent's} preference weights. \\
\textit{Preference weights / utility} & In the strategic ledger, declared weights $w_{i,\mathrm{cash}},w_{i,k}$ map cash and metrics to utility $U_{i,t}$; utility need not equal profit. \\
\textit{Resource / constraint} & An available balance (such as cash $b_{i,t}$) and an enforced limit excluding unaffordable actions. \\
\bottomrule
\end{tabularx}
\end{table}

For instance, consider a partner deciding whether to fund a pilot. The
partner is an actor whose goal is preserving capital while gaining credible
participation; its preference weights quantify that trade-off, whereas action
costs and credibility gains create incentives. Funding and waiting are actions,
and selecting one from the actor's observation is a decision whose consequences
the ledger settles. This example illustrates how the terms connect rather than
claiming empirical partner behavior.

Moreover, within the same cast, actors can assign different preference
weights to retained cash, participation, credibility, or a declared risk
metric. The author must state what each metric represents and why its weight
reflects an assumed trade-off, because calling a persona ``cautious'' in prose
does not enforce caution in settlement. In everyday terms, \textit{utility}
resembles commuters ranking routes: one commuter weights a lower fare more
heavily, whereas another weights a shorter trip, so different choices emerge
without assuming everyone minimizes fare.

Beyond individual preference weights, strategic dependence requires rules
that couple actors' choices. Supporting a launch may pay off only if an
independent reviewer verifies it, or two simultaneous actions may shift a
shared metric more than the sum of their individual effects. However, a
promise or threat in persona text has no binding force unless an executable
rule enforces it. Similarly, splitting one actor cohort into multiple accounts
requires partitioning its balances and per-action effects rather than copying
its full market power into each new account.

\subsection{\rev{Social Graph}, Information, Rules, and Event Timing}
Alongside incentives, each \rev{agent in the swarm} operates behind an explicit
observation boundary \rev{defined by a \textit{social graph}}. Directed edges in
the \rev{\textit{social graph} (implemented as an \textit{information graph})
determine how agents in the swarm communicate and whose messages an agent
receives}, while the engine withholds other \rev{agents'} private preference
weights and unannounced future events. In the strategic ledger, messages
attached to round-$t$ actions arrive in round $t+1$, ensuring that query order
within a round never grants an unintended information advantage. By analogy, an
\textit{observation} is a card player's view of their own hand and public table
cards, whereas \rev{the swarm's \textit{social graph}} works like mailing-list
subscriptions, which need neither be mutual nor imply trust.
Table~\ref{tab:decision-definitions} summarizes these terms.

\begin{table}[H]
\caption{\rev{Decisions, social-graph observations, and rules.}}
\label{tab:decision-definitions}
\centering\small
\renewcommand{\arraystretch}{1.12}
\begin{tabularx}{\columnwidth}{@{}>{\raggedright\arraybackslash}p{0.25\columnwidth}Y@{}}
\toprule
Term & Definition \\
\midrule
\textit{Action} & A permitted option $a \in \operatorname{Menu}_i$ with declared cost $c_i(a)$, effects, and optional preset message; \texttt{wait} is an action. \\
\textit{Decision} & Selection of an affordable action $a_{i,t} \in A_{i,t}$ from \rev{agent} $i$'s observation; the ledger settles consequences separately. \\
\textit{Observation} & Information visible to one \rev{agent} at $t$: its own profile and balances, public state, visible events, and delivered \rev{social-graph} messages. \\
\rev{\textit{Social / information graph}} & \rev{Directed communication network through which the \textit{swarm of agents} exchanges messages, specifying permitted delivery rather than measuring social influence.} \\
\textit{Event} & An exogenous change with declared timing, audience visibility, and state effects. \\
\textit{Rule / interaction} & A settlement law; in the strategic ledger, a joint-action rule $r$ adds effect $g_{k,r}$ when a specified action combination occurs. \\
\bottomrule
\end{tabularx}
\end{table}

To settle choices consistently, the world also declares action costs,
metric bounds $[L_k,H_k]$, ledger parameters, and exogenous event schedules.
Default strategic worlds apply deterministic event schedules, whereas
probabilistic event trees belong to a separate event-chain module. Keeping
exogenous events distinct from endogenous actor actions lets later analysis
trace whether a trajectory shift came from an external shock or from cast
reactions.

\subsection{Evidence, Assumptions, and Uncertainty}
Because settlement rules fix the meaning of every simulated number, each
material parameter must trace back to cited evidence or a labeled assumption.
An \textit{assumption} is an unverified modeling input, whereas
\textit{evidence} is a dated external record motivating part of the world. For
example, a parliamentary agenda may establish the date of a debate without
establishing how traders will react; world construction therefore records the
dated source and the assumed behavioral response separately, following
established model-documentation practice \cite{grimm2020}.

On the other hand, whenever a scenario introduces numerical probabilities,
the author must distinguish a user-chosen scenario branch, a model-returned
option score $q_{i,t}(a)$, a seeded sampling draw, and an empirical frequency.
In particular, \textit{Bayesian} updating requires declared hypotheses, priors,
and likelihoods \cite{pishronik2014} rather than inferring odds from narrative
text; uncalibrated scenario weights are therefore reported strictly as
assumptions.

\subsection{Construction Review and Output}
Before any run starts, automated checks and human review audit the assembled
world. Specifically, construction checks verify balance conservation, metric
units, affordable action menus, isolated unilateral deltas, information links,
and joint-action couplings, flagging degenerate setups such as actors with only
one affordable option or goals unaffected by any permitted action.

Consequently, passing construction review guarantees an explicit,
executable specification rather than an empirically calibrated population;
human judgment remains responsible for checking whether preference weights match
intended roles and whether cited sources justify the parameters. With that
specification frozen, Section~\ref{sec:method-system} details how
\textit{Boyun} executes rounds and decision points locally.
